\documentclass[memorial,capa]{texifsap}

\usepackage[utf8]{inputenc} % acentuacao
\usepackage{graphicx} % para inserir figuras
\usepackage{array} % células de quadros alinhadas à esquerda: >{\raggedright\arraybackslash}p{...}
\graphicspath{{imagens/}}
\usepackage[T1]{fontenc}
% coluna de largura fixa com texto alinhado à esquerda (evita espaços largos)
\newcolumntype{L}[1]{>{\raggedright\arraybackslash}p{#1}}

%==============================================================================
% MEMORIAL DESCRITIVO DE VALIDAÇÃO DE CONHECIMENTOS E EXPERIÊNCIAS
% PROFISSIONAIS — APROVEITAMENTO DO TRABALHO DE CONCLUSÃO DE CURSO
%
% REGRAS (Regulamento do TCC, Capítulo VII, arts. 47 a 52):
% - O memorial descreve, de forma DETALHADA, as atividades profissionais
%   e/ou cursos de educação profissional e tecnológica relacionados à área
%   de Análise e Desenvolvimento de Sistemas, alinhados ao perfil do
%   egresso do curso (PPC) e ao programa das disciplinas (art. 49, I).
% - Valem experiências de, no MÍNIMO, DOIS ANOS, ocorridas nos ÚLTIMOS DEZ
%   ANOS, contados da data da solicitação (art. 49, II).
% - Texto acadêmico-científico: relato das atividades COM FUNDAMENTAÇÃO
%   TEÓRICA, com o período de cada atividade e as devidas comprovações;
%   ENTRE CINCO E DEZ PÁGINAS (art. 49, III).
% - Comprovantes OBRIGATÓRIOS, anexados ao memorial: Carteira de Trabalho
%   (CTPS), carta da chefia/superior ou documento que comprove atividade
%   autônoma (carta de contador e/ou documento da Receita Federal)
%   (art. 49, IV).
% - O pedido é feito por requerimento à Coordenação de Registros
%   Acadêmicos, durante o período letivo, estando regularmente matriculado;
%   só é aceito UM ÚNICO pedido (art. 48, §§ 1º e 3º; art. 49, V).
% - Uma comissão de três docentes avalia o memorial e pode exigir
%   apresentação oral de 20 a 25 minutos, com arguição (arts. 50 e 51).
%
% O QUE OS MEMORIAIS BEM AVALIADOS TÊM EM COMUM (detalhes e exemplos em
% instruções/como_elaborar_um_bom_memorial.txt):
%  1. período EXATO (mês/ano) de cada cargo, com a progressão na carreira;
%  2. contexto da organização e do sistema (arquitetura, tecnologias, equipe);
%  3. atividades descritas como DEMANDA -> AÇÃO -> RESULTADO, com o impacto
%     quantificado sempre que possível;
%  4. FUNDAMENTAÇÃO em autores clássicos da área, no padrão
%     "Segundo Autor (ano), <conceito>. Na minha atuação, <prática>.";
%  5. DISCIPLINAS do curso citadas pelo nome, junto às atividades;
%  6. EVIDÊNCIAS concretas: quadros com artefatos anonimizados (consulta SQL,
%     cenário de teste, etapas de projeto) e links de repositórios públicos;
%  7. competências listadas ao final de cada experiência;
%  8. sigilo da organização preservado (dados de clientes, código, valores);
%  9. conclusão que sintetiza a trajetória e formaliza o pedido.
%==============================================================================

\unidade{C{\^a}mpus Sapucaia do Sul}

%TÍTULO DE FORMAÇÃO: descomente o bloco correspondente ao seu curso e
%mantenha o outro comentado.
%
%--- Curso Superior de Tecnologia em Análise e Desenvolvimento de Sistemas:
\curso{An\'alise e Desenvolvimento de Sistemas}
\nomecurso{Superior de Tecnologia em An\'alise e Desenvolvimento de Sistemas}
%na capa, o nome do curso quebra a linha no ponto indicado pelo \\
\nomecursocapa{Superior de Tecnologia em\\An\'alise e Desenvolvimento de Sistemas}
\titulocurso{Tecn\'ologo(a) em An\'alise e Desenvolvimento de Sistemas}

%TÍTULO: mantenha como está. A segunda linha (depois do \\) é o subtítulo.
\title{Memorial descritivo de valida\c{c}\~ao de conhecimentos e experi\^encias profissionais\\[.5cm]Aproveitamento do Trabalho de Conclus\~ao de Curso}

%AUTOR: o memorial é individual e não tem orientador.
\author{ÚltimoSobrenome}{Primeiro nome e sobrenome}

\begin{document}

\maketitle

\sloppy

%Sumario
\tableofcontents

%OPCIONAL: se usar muitos quadros, tabelas ou figuras, descomente a lista
%correspondente ANTES do \tableofcontents (as listas vêm antes do sumário).
%\listofquadros
%\listoftables
%\listoffigures

%==============================================================================
% PARTE TEXTUAL — entre CINCO e DEZ páginas (Regulamento do TCC, art. 49,
% III). Organização recomendada:
%   1 Introdução
%   2 Trajetória profissional (quadro-resumo com períodos e comprovantes)
%   3, 4, ... Um capítulo por experiência profissional (ou curso de
%     educação profissional e tecnológica), em ordem cronológica
%   n Relação com o perfil do egresso e com as disciplinas do curso
%   n+1 Considerações finais (com a solicitação de aproveitamento)
%   Referências e Anexos (comprovantes)
% Outras organizações também são bem avaliadas: (a) um capítulo de
% fundamentação teórica aplicada antes da trajetória; ou (b) um capítulo de
% contexto profissional seguido de um capítulo POR ÁREA DE CONHECIMENTO
% (requisitos, gestão de projetos, APIs, banco de dados, segurança...),
% cada uma com teoria, prática e disciplinas. Escolha uma e mantenha-a.
% Pessoa do discurso: 1ª pessoa ("atuei") ou 3ª ("o autor atuou") são
% aceitas; apenas não as misture.
%==============================================================================

%ATENÇÃO — PLÁGIO E INTELIGÊNCIA ARTIFICIAL (Regulamento do TCC, art. 34):
%plágio total ou parcial ANULA o TCC e reprova o acadêmico. O uso de
%ferramentas de inteligência artificial para a produção do trabalho é
%considerado plágio. Escreva o texto com suas próprias palavras e cite e
%referencie TODAS as fontes utilizadas.

\chapter{Introdução}

%GUIA: apresente o objetivo do memorial e a sua base regulamentar
%(Capítulo VII do Regulamento do TCC), resuma a sua formação prévia, se
%for relevante (outros cursos, graduações, certificações), e dê uma visão
%geral da trajetória que será descrita: quais experiências, em que
%períodos e em que áreas. No último parágrafo, descreva a organização do
%texto.

O Trabalho de Conclusão de Curso (TCC) é atividade curricular
obrigatória do Curso Superior de Tecnologia em Análise e Desenvolvimento
de Sistemas do Instituto Federal Sul-rio-grandense (IFSul), Câmpus
Sapucaia do Sul. O Capítulo VII do Regulamento do Trabalho de Conclusão
de Curso~\cite{IFSul:RegTCC}, que integra o Projeto Pedagógico do
Curso~\cite{IFSul:PPC}, prevê, contudo, que os conhecimentos
adquiridos na experiência profissional possam ser avaliados,
reconhecidos e certificados para fins de aproveitamento do TCC. Nesse
contexto, este memorial descritivo tem por objetivo apresentar e
fundamentar as experiências profissionais do autor na área de Análise e
Desenvolvimento de Sistemas, demonstrando a sua correspondência com as
competências previstas no perfil do egresso do curso.

Descreva aqui, em um ou dois parágrafos, a sua formação e trajetória de
forma resumida: há quanto tempo atua na área, em quais organizações, em
que funções e com quais tecnologias e metodologias. Se possuir formação
anterior (outros cursos técnicos ou superiores, especializações,
certificações) que ajude a contextualizar a trajetória, mencione-a aqui.

O memorial está organizado da seguinte forma: o Capítulo~\ref{cap:trajetoria}
apresenta a trajetória profissional e os períodos de cada experiência;
os Capítulos~\ref{cap:experiencia1} e~\ref{cap:experiencia2} descrevem as
atividades desenvolvidas em cada organização, com a respectiva
fundamentação teórica; o Capítulo~\ref{cap:egresso} relaciona as
competências desenvolvidas ao perfil do egresso e às disciplinas do
curso; e o Capítulo~\ref{cap:consideracoes} apresenta as considerações
finais e a solicitação de aproveitamento.

\chapter{Trajetória profissional}
\label{cap:trajetoria}

%GUIA: apresente uma visão geral da trajetória, em ordem cronológica, e
%resuma TODAS as experiências no quadro abaixo, com o período exato
%(mês/ano de início e de fim) e o documento comprobatório de cada uma,
%identificado pela letra do anexo correspondente. O quadro permite à
%comissão verificar rapidamente os requisitos do art. 49, II: período
%mínimo de DOIS ANOS de experiência, ocorrida nos ÚLTIMOS DEZ ANOS.
%Cursos de educação profissional e tecnológica também podem ser
%incluídos (art. 49, I), indicando a carga horária e o certificado.
%Remova do quadro as linhas que não se aplicam (atividade autônoma, curso)
%e os anexos correspondentes.

Apresente aqui, em linhas gerais, a evolução da sua trajetória na área
de tecnologia: como ingressou na área, como as funções e as
responsabilidades evoluíram e quais foram os principais domínios de
atuação (por exemplo, desenvolvimento de software, banco de dados,
testes e qualidade, infraestrutura, análise de requisitos ou gestão de
projetos).

O Quadro~\ref{quadro-trajetoria} resume as experiências descritas neste
memorial, com os respectivos períodos e documentos comprobatórios,
reunidos nos anexos. Somadas, as experiências totalizam N anos e M
meses de atuação na área, todos nos últimos dez anos.

\begin{quadro}[htbp]
\caption{Experiências profissionais descritas neste memorial}
\label{quadro-trajetoria}
\centering\begin{tabular}{|L{3.2cm}|L{2.8cm}|L{2.8cm}|L{1.9cm}|L{2.9cm}|}
\hline
\textbf{Organização ou curso} & \textbf{Cargo/função} & \textbf{Período} & \textbf{Duração} & \textbf{Comprovante}\\
\hline
{\small Nome da organização 1} & {\small Cargo exercido} & {\small mm/aaaa a mm/aaaa} & {\small N meses} & {\small Anexo~\ref{anx:ctps}}\\
\hline
{\small Nome da organização 2} & {\small Cargo exercido} & {\small mm/aaaa até o presente} & {\small N meses} & {\small Anexos~\ref{anx:ctps2} e~\ref{anx:carta}}\\
\hline
{\small Atividade autônoma (se houver)} & {\small Serviço prestado} & {\small mm/aaaa a mm/aaaa} & {\small N meses} & {\small Anexo~\ref{anx:contador}}\\
\hline
{\small Curso de qualificação profissional (se houver)} & {\small Nome do curso (N horas)} & {\small mm/aaaa a mm/aaaa} & {\small --} & {\small Anexo~\ref{anx:certificado}}\\
\hline
\multicolumn{3}{|r|}{\small\textbf{Total}} & {\small\textbf{N meses}} & \\
\hline
\end{tabular}
\begin{flushleft}
{\small Fonte: Elaboração própria.}
\end{flushleft}
\end{quadro}

%==============================================================================
% UM CAPÍTULO POR EXPERIÊNCIA. Copie o capítulo abaixo para cada
% organização (ou curso), em ordem cronológica. Se houve mudança de cargo
% na mesma organização, use uma seção para cada cargo, com o respectivo
% período.
%==============================================================================

\chapter{Experiência profissional em Nome da organização 1}
\label{cap:experiencia1}

%GUIA: indique no primeiro parágrafo o PERÍODO (mês/ano de início e de
%fim), o CARGO e a carga horária semanal, e o documento comprobatório.
%Depois, contextualize a organização e o sistema ou produto em que atuou.
%Não exponha informações sigilosas da organização (dados de clientes,
%código-fonte proprietário, valores, credenciais).

Entre mm/aaaa e mm/aaaa, o autor atuou na Nome da organização 1, no
cargo de Cargo exercido, com jornada de N horas semanais, conforme a
Carteira de Trabalho apresentada no Anexo~\ref{anx:ctps}. Apresente aqui
a organização: o seu ramo de atuação, o porte, os clientes atendidos e o
papel da área de tecnologia no negócio.

\section{Contexto técnico}
\label{sec:contexto1}

%GUIA: descreva o sistema, produto ou ambiente em que atuou: a
%arquitetura, as linguagens, os frameworks, os bancos de dados, a
%infraestrutura e as ferramentas, e a organização da equipe (tamanho,
%papéis, metodologia de trabalho).

Descreva aqui o sistema ou o ambiente em que atuou, com a arquitetura,
as tecnologias empregadas e a organização da equipe. Por exemplo, se a
equipe adotava o Scrum, explique brevemente o framework e o seu papel
nele: segundo \citet{Schwaber:2020}, o Scrum organiza o trabalho em
ciclos de duração fixa, as \textit{sprints}, nas quais um time
multidisciplinar entrega incrementos do produto e inspeciona
continuamente o progresso.

\section{Atividades desenvolvidas}
\label{sec:atividades1}

%GUIA: esta é a parte central do memorial. Para cada atividade
%relevante, descreva: (i) o problema ou a demanda; (ii) o que você fez,
%com quais técnicas e ferramentas; (iii) o resultado obtido; e (iv) a
%FUNDAMENTAÇÃO TEÓRICA — os conceitos da área que sustentam a atividade,
%com citação de autores (livros, artigos, normas). A fundamentação é
%exigência do Regulamento (art. 49, III) e o que diferencia o memorial
%de um currículo. Use subseções se houver muitas atividades.
%
%PADRÃO QUE FUNCIONA, para cada atividade:
%  "Segundo Autor (ano), <conceito da área>. Na minha atuação, <o que
%   foi feito, com quais ferramentas>. Como resultado, <impacto, de
%   preferência com números: tempo, volume, falhas, equipe>. Essa
%   atividade relaciona-se à disciplina <nome da disciplina>."
%Mostre EVIDÊNCIAS: um quadro com um artefato real anonimizado (consulta
%SQL, cenário BDD, etapas do projeto, trecho de configuração) e links de
%repositórios públicos, quando houver. Explique em nota de rodapé
%(\footnote{...}) termos próprios da organização ou do produto.
%Evite: listar tecnologias sem dizer o que fez com elas; frases genéricas
%("trabalhei com várias tecnologias"); expor dados sigilosos.
%
%Citações (NBR 10520:2023): \cite{chave} -> (Sommerville, 2019);
%\citet{chave} -> Sommerville (2019); com página: \cite[p.~45]{chave}.

Descreva aqui cada atividade relevante, relacionando-a aos conceitos da
área. Por exemplo, ao relatar o levantamento de requisitos de uma nova
funcionalidade, explique que a engenharia de requisitos compreende a
elicitação, a análise, a especificação e a validação das necessidades
dos usuários~\cite{Sommerville:2019}, e mostre como cada etapa ocorreu na
sua prática.

Da mesma forma, ao descrever o desenvolvimento ou o teste de uma API,
fundamente-o no estilo arquitetural adotado. \citet{Fielding:2000}
definiu o estilo REST por um conjunto de restrições, como a comunicação
cliente-servidor sem estado e a interface uniforme, que orientam o
projeto de serviços web. Relate então o que foi implementado e quais
resultados foram alcançados, como a redução de falhas em produção ou o
ganho de desempenho, quantificando-os sempre que possível, e indique a
disciplina do curso que trata do mesmo conteúdo (por exemplo,
Desenvolvimento de Sistemas Computacionais).

\section{Competências desenvolvidas}

%GUIA: liste, em alíneas, as competências técnicas e comportamentais
%desenvolvidas nessa experiência. Elas serão retomadas no quadro do
%capítulo de relação com o perfil do egresso.

A atuação na Nome da organização 1 contribuiu para o desenvolvimento das
seguintes competências:
\begin{itemize}
\item competência técnica 1 (por exemplo, desenvolvimento de aplicações
  web com a tecnologia X);
\item competência técnica 2 (por exemplo, modelagem e consulta de bancos
  de dados relacionais);
\item competência comportamental (por exemplo, trabalho em equipe com
  metodologias ágeis).
\end{itemize}

\chapter{Experiência profissional em Nome da organização 2}
\label{cap:experiencia2}

Entre mm/aaaa e o presente, o autor atua na Nome da organização 2, no
cargo de Cargo exercido, conforme a Carteira de Trabalho e a carta da
chefia imediata apresentadas nos Anexos~\ref{anx:ctps2}
e~\ref{anx:carta}. Apresente aqui a organização e o seu contexto, como
no capítulo anterior.

\section{Contexto técnico}

Descreva aqui o sistema ou o ambiente em que atua, as tecnologias e a
organização da equipe. Se houver práticas de garantia da qualidade,
relacione-as à literatura: para \citet{Pressman:2021}, a qualidade de
software envolve a conformidade com os requisitos, com padrões de
desenvolvimento e com as características esperadas de um produto
profissional.

\section{Atividades desenvolvidas}
\label{sec:atividades2}

Descreva aqui as atividades desenvolvidas, com a mesma estrutura do
capítulo anterior: a demanda, o que foi feito, o resultado e a
fundamentação teórica. Se houve progressão de cargo na organização,
crie uma seção para cada cargo, indicando o respectivo período.

Sempre que possível, apresente uma evidência concreta da atividade. Por
exemplo, ao relatar a validação de dados de um sistema, explique que o
modelo relacional organiza os dados em tabelas relacionadas entre si,
cuja consistência é garantida por restrições de
integridade~\cite{Elmasri:2018}, e mostre uma consulta utilizada, com os
nomes de tabelas e os identificadores anonimizados, como no
Quadro~\ref{quadro-evidencia}. Em seguida, explique a finalidade da
consulta e o resultado obtido com ela.

\begin{quadro}[htbp]
\caption{Exemplo de consulta SQL utilizada na validação de dados (nomes anonimizados)}
\label{quadro-evidencia}
\centering\begin{tabular}{|p{12cm}|}
\hline
{\ttfamily\small SELECT p.id, p.situacao, u.nome AS responsavel}\\
{\ttfamily\small FROM TABELA\_PEDIDOS p}\\
{\ttfamily\small JOIN TABELA\_USUARIOS u ON u.id = p.id\_usuario}\\
{\ttfamily\small WHERE p.situacao = 'PENDENTE'}\\
{\ttfamily\small \ \ AND p.data\_criacao < CURRENT\_DATE - 30;}\\
\hline
\end{tabular}
\begin{flushleft}
{\small Fonte: Elaboração própria.}
\end{flushleft}
\end{quadro}

\section{Competências desenvolvidas}

A atuação na Nome da organização 2 tem contribuído para o
desenvolvimento das seguintes competências:
\begin{itemize}
\item competência técnica 1;
\item competência técnica 2;
\item competência comportamental.
\end{itemize}

\chapter{Relação com o perfil do egresso e com as disciplinas do curso}
\label{cap:egresso}

%GUIA: este capítulo demonstra o requisito central da validação: a
%correspondência entre as experiências e as habilidades e competências
%do PERFIL DO EGRESSO, descritas no Projeto Pedagógico do Curso (PPC)
%(art. 48), e o programa das disciplinas (art. 49, I). Inclua no quadro
%as competências que as suas atividades EVIDENCIAM, transcritas do PPC
%sem reescrevê-las, e indique, para cada uma, as atividades dos capítulos
%anteriores e as disciplinas relacionadas. Uma competência pode ser
%evidenciada por várias atividades e relacionar-se a várias disciplinas.
%
%PERFIL DO EGRESSO (PPC 2023, seção 3.6, p. 41-42), para consulta:
%"O perfil profissional do(a) egresso(a) do Curso contempla o domínio de
%trabalho em equipes com iniciativa, criatividade e sociabilidade,
%realizando suas atividades profissionais de forma ética, com
%responsabilidade socioambiental, atendendo às normas técnicas e de
%segurança, compreendendo:
% - habilidade para aprender e lidar com o inusitado, conviver e se
%   comunicar;
% - capacidade para aprender de forma autônoma e contínua;
% - comprometimento com a melhoria da qualidade de vida;
% - conhecimento sobre importância de agir com ética, comprometimento
%   socioambiental e solidariedade enquanto ser humano, cidadão e
%   profissional;
% - o gerenciamento de processos participativos de organização pública
%   e/ou privada e/ou incluir-se neles;
% - conhecimento sobre a importância de agir com responsabilidade,
%   sensibilidade e equilíbrio ao enfrentar desafios profissionais.
%Na atuação deste(a) profissional, destacam-se as seguintes atividades:
% - analisar, projetar, desenvolver, testar, implantar e manter sistemas
%   computacionais de informação;
% - avaliar, selecionar, especificar e utilizar metodologias, tecnologias
%   e ferramentas da Engenharia de Software, linguagens de programação e
%   bancos de dados;
% - coordenar equipes de produção de softwares;
% - vistoriar, realizar perícia, avaliar, emitir laudo e parecer técnico
%   em sua área de formação."
%DISCIPLINAS DO CURSO POR ÁREA (Matriz curricular nº 2, a partir de
%2023/2, e matriz de eletivas — PPC, Apêndices D e E). Use os nomes
%exatamente como abaixo; (E) = eletiva:
% - Programação e algoritmos: Computação Aplicada à Solução de Problemas;
%   Princípios do Pensamento Computacional; Programação Orientada a
%   Objetos; Estruturas de Dados; Análise de Algoritmos (E).
% - Desenvolvimento de sistemas e interfaces: Desenvolvimento de
%   Interfaces I e II; Fundamentos de Desenvolvimento de Sistemas
%   Computacionais; Desenvolvimento de Sistemas Computacionais;
%   Desenvolvimento de aplicações para Dispositivos Móveis; Tópicos
%   Avançados em Programação para Dispositivos Móveis (E); Experiência do
%   Usuário e Acessibilidade (E).
% - Engenharia de software, qualidade e projetos: Fundamentos de
%   Engenharia de Sistemas Computacionais; Engenharia de Sistemas
%   Computacionais; Qualidade de Sistemas Computacionais; Gestão de
%   Projetos; Projeto de Sistemas Computacionais (E).
% - Banco de dados e análise de dados: Fundamentos de Sistemas de Banco de
%   Dados; Tópicos Avançados em Sistemas de Banco de Dados; Ciência e
%   Análise de Dados.
% - Infraestrutura, redes e segurança: Fundamentos Computacionais;
%   Fundamentos de Infraestrutura Computacional; Fundamentos de Segurança
%   de Sistemas Computacionais; Segurança de Sistemas Computacionais;
%   Tópicos Avançados em Segurança de Sistemas Computacionais (E);
%   Internet das Coisas (E).
% - Inteligência artificial: Inteligência Artificial (E).
% - Gestão, negócios e relações de trabalho: Relações Interpessoais;
%   Empreendedorismo e Inovação; Gestão Econômico-Financeira; Economia;
%   Custos Empresariais (E); Segurança no Trabalho.
% - Formação geral: Fundamentos de Matemática Discreta; Matemática
%   Vetorial e Matricial; Metodologia Científica e de Extensão; Leitura e
%   Produção Textual; Introdução à Ética; Computação e Sociedade I e II.
%(Seminários I, Seminários II e o TCC são os componentes que estão sendo
%validados; não os use no quadro.)
%Campo de atuação: empresas de planejamento, desenvolvimento de projetos,
%assistência técnica e consultoria, empresas de tecnologia, empresas em
%geral (indústria, comércio e serviços), ONGs, órgãos públicos,
%institutos e centros de pesquisa e instituições de ensino.

O perfil do egresso do Curso Superior de Tecnologia em Análise e
Desenvolvimento de Sistemas, descrito no Projeto Pedagógico do
Curso~\cite[p.~41-42]{IFSul:PPC}, contempla o domínio de trabalho em
equipes, realizado de forma ética e atendendo às normas técnicas, e
destaca, na atuação do profissional, atividades como analisar,
projetar, desenvolver, testar, implantar e manter sistemas
computacionais de informação. O Quadro~\ref{quadro-egresso} relaciona essas
competências às atividades descritas neste memorial e às disciplinas do
curso que tratam dos mesmos conteúdos.

\begin{quadro}[htbp]
\caption{Relação entre as competências do perfil do egresso, as atividades profissionais e as disciplinas do curso}
\label{quadro-egresso}
\centering\begin{tabular}{|L{5.4cm}|L{4.8cm}|L{4.2cm}|}
\hline
\textbf{Competência do perfil do egresso (PPC)} & \textbf{Atividades que a evidenciam} & \textbf{Disciplinas relacionadas}\\
\hline
{\small Analisar, projetar, desenvolver, testar, implantar e manter sistemas computacionais de informação} & {\small Atividades da Seção~\ref{sec:atividades1} (levantamento de requisitos e desenvolvimento da API)} & {\small Fundamentos de Engenharia de Sistemas Computacionais; Desenvolvimento de Sistemas Computacionais}\\
\hline
{\small Avaliar, selecionar, especificar e utilizar metodologias, tecnologias e ferramentas da Engenharia de Software, linguagens de programação e bancos de dados} & {\small Atividade da Seção~\ref{sec:atividades2} (validação de dados)} & {\small Tópicos Avançados em Sistemas de Banco de Dados}\\
\hline
{\small Coordenar equipes de produção de softwares} & {\small Descreva a atividade (por exemplo, liderança técnica de uma equipe)} & {\small Gestão de Projetos}\\
\hline
{\small Domínio de trabalho em equipes com iniciativa, criatividade e sociabilidade} & {\small Atuação em equipe com metodologia ágil (Seção~\ref{sec:contexto1})} & {\small Relações Interpessoais}\\
\hline
\end{tabular}
\begin{flushleft}
{\small Fonte: Elaboração própria, com base em \citet{IFSul:PPC}.}
\end{flushleft}
\end{quadro}

Comente aqui o quadro, destacando as competências mais desenvolvidas na
trajetória e a forma como a experiência profissional complementou os
conteúdos estudados no curso.

\chapter{Considerações finais}
\label{cap:consideracoes}

%GUIA: sintetize a trajetória descrita, retome as principais competências
%desenvolvidas e mostre que elas correspondem às esperadas para o TCC.
%Vale destacar como a experiência profissional e o curso se
%complementaram e reconhecimentos formais recebidos (promoções,
%premiações). Não apresente experiências ou dados novos que não tenham
%aparecido antes.

Sintetize aqui a trajetória descrita no memorial, em ordem cronológica,
e as principais competências desenvolvidas em cada experiência,
evidenciando que, em conjunto, elas abrangem as dimensões do
desenvolvimento de software esperadas do egresso, como a análise de
requisitos, a modelagem de dados, o projeto e a implementação de
sistemas, os testes, a manutenção e o trabalho em equipe.

\section{Solicitação de aproveitamento}

Diante do exposto, e considerando que as experiências profissionais
documentadas neste memorial correspondem às habilidades e competências
previstas no perfil do egresso do Curso Superior de Tecnologia em
Análise e Desenvolvimento de Sistemas, o autor solicita o aproveitamento
do Trabalho de Conclusão de Curso por meio da validação de
conhecimentos e experiências profissionais, na forma prevista no
Capítulo VII do Regulamento do Trabalho de Conclusão de Curso do Curso
Superior de Tecnologia em Análise e Desenvolvimento de Sistemas do
Instituto Federal Sul-rio-grandense, Câmpus Sapucaia do Sul.

\section{Declaração sobre uso de Inteligência Artificial}

%GUIA (Portaria CNPq nº 2.664/2026): declare o uso de ferramentas de
%inteligência artificial generativa em qualquer fase da elaboração do
%memorial, com a ferramenta, a finalidade e a fase. Lembre-se do
%Regulamento do TCC (art. 34): o uso de IA para a PRODUÇÃO do trabalho é
%considerado plágio. Veja instruções/como_declarar_uso_de_ia.txt.

Declare aqui, de forma transparente, o uso de ferramentas de
inteligência artificial generativa na elaboração deste memorial,
especificando a ferramenta, com nome e versão, a finalidade e a fase do
trabalho em que ocorreu. Caso nenhuma ferramenta tenha sido utilizada,
declare-o expressamente.

%REFERÊNCIAS: liste somente as obras citadas no texto (NBR 6023). As
%regras de preenchimento do .bib estão no final de tcc/exemplo-tcc.tex.
\bibliographystyle{abnt}
\bibliography{bibliografia-memorial}

%==============================================================================
% ANEXOS — DOCUMENTOS COMPROBATÓRIOS (OBRIGATÓRIOS; Regulamento do TCC,
% art. 49, IV): Carteira de Trabalho (páginas de identificação e do
% contrato de trabalho), carta da chefia/superior e/ou documento que
% comprove atividade autônoma (carta de contador, documento da Receita
% Federal). Cursos de educação profissional e tecnológica: certificado
% com carga horária. Um anexo para cada documento, na ordem em que são
% citados no texto.
%
% Para inserir um documento em PDF, use (uma página por vez ou todas):
%   \includepdf[pages=-,pagecommand={}]{anexos/ctps.pdf}
% (o pacote pdfpages já é carregado pela classe). Se preferir imagens:
%   \includegraphics[width=\textwidth]{anexos/ctps.png}
% ATENÇÃO À PRIVACIDADE: oculte dados pessoais que não sejam necessários à
% comprovação (como CPF, número de documentos, endereço e salário).
%==============================================================================
\anexos

%Cada anexo recebe a cópia digitalizada de um documento comprobatório
%previsto no Regulamento do TCC (art. 49, IV) ou no PPC (seção 3.8.12).
%Substitua o texto de cada anexo pelo documento real e remova os anexos
%que não se aplicam (e as linhas correspondentes do Quadro 1).

\chapter{Carteira de Trabalho -- Nome da organização 1}
\label{anx:ctps}

Insira aqui a cópia da Carteira de Trabalho e Previdência Social (CTPS)
que comprova o período de atuação na Nome da organização 1, gerada pelo
aplicativo Carteira de Trabalho Digital ou digitalizada da carteira em
papel (página de identificação e página do contrato de trabalho).

\chapter{Carteira de Trabalho -- Nome da organização 2}
\label{anx:ctps2}

Insira aqui a cópia da Carteira de Trabalho que comprova o período de
atuação na Nome da organização 2.

\chapter{Carta da chefia -- Nome da organização 2}
\label{anx:carta}

Insira aqui a carta da chefia imediata ou de superior da organização,
assinada, que declare o cargo, o período e as principais atividades
desempenhadas pelo autor.

\chapter{Declaração de contador -- atividade autônoma}
\label{anx:contador}

Se houver atividade autônoma, insira aqui a carta do contador e/ou o
documento da Receita Federal que comprove a atividade e o seu período.

\chapter{Certificado de curso -- Nome do curso}
\label{anx:certificado}

Se o memorial incluir curso de educação profissional e tecnológica,
insira aqui a cópia do certificado, com a carga horária e o período de
realização.

\end{document}
